% GENERATED FILE. Edit canonical lessons, then run npm run generate:books.
% canonical-source-hash: fnv1a64:60d61ed48c0008a0
% canonical-lessons: BN-C266-porchhilam, BN-C266-jokhon-tokhon, BN-C266-thik, BN-W266-letter-ttha, BN-C266-otha, BN-R266-the-scene

\chapter{What Was Going On}
\label{ch:bn-what-was-going-on}
\hlchaptermodality{266}

\begin{chapteropening}
\textbf{By the end of this chapter:} \emph{I can say what is going on and what was going on (pôṛchhi, pôṛchhilām), set a scene with jôkhon … tôkhon, say oṭhā, to get up, and read and write \bn{ঠ} in \bn{ঠিক}.}

Complete BN-R266-the-scene: say what was going on, and what broke into it.
\end{chapteropening}

\section[\texorpdfstring{\bn{আমি} \bn{পড়ছিলাম} (āmi pôṛchhilām)}{আমি পড়ছিলাম (āmi pôṛchhilām)}]{\bn{আমি} \bn{পড়ছিলাম} (āmi pôṛchhilām) — going on now, going on then}

\label{lesson:BN-C266-porchhilam}



\begin{quote}
Say \textquotedblleft{}I was at home\textquotedblright{} and \textquotedblleft{}I will go\textquotedblright{}. (\emph{chhilām}, \emph{jābô}.) Both are whole facts. This lesson is about things in the middle of happening.
\end{quote}

\subsection*{You'll want to know: going on now}
\begin{quote}
\textbf{\bn{আমি} \bn{এখন} \bn{একটা} \bn{বই} \bn{পড়ছি।}}
\end{quote}

\begin{quote}
\emph{āmi ekhôn ekṭā bôi pôṛchhi.} — \textquotedblleft{}I'm reading a book now.\textquotedblright{}
\end{quote}

What is going on at this moment takes \textbf{-\bn{ছ}-} after the stem: \textbf{\bn{আমি} \bn{পড়ছি}} (\emph{āmi pôṛchhi}), \textbf{\bn{তুমি} \bn{পড়ছ}} (\emph{tumi pôṛchhô}), \textbf{\bn{সে} \bn{পড়ছে}} (\emph{she pôṛchhe}), \textbf{\bn{আপনি} \bn{পড়ছেন}} (\emph{āpni pôṛchhen}). After a stem that ends in a vowel it is \textbf{-\bn{চ্ছ}-}: \textbf{\bn{যাচ্ছি}} (\emph{jāchchhi}), \textquotedblleft{}I'm going\textquotedblright{}; \textbf{\bn{সে} \bn{খাচ্ছে}} (\emph{she khāchchhe}); \textbf{\bn{কী} \bn{হচ্ছে}?} (\emph{kī hôchchhe?}), \textquotedblleft{}what's happening?\textquotedblright{}. An \textbf{\bn{এ}} said \emph{e} rises, as it does in the past: \textbf{\bn{শেখা}} $\to$ \textbf{\bn{শিখছি}} (\emph{shikhchhi}).

\subsection*{You'll want to know: going on then}
Put \textbf{-\bn{ছিল}-} where \textbf{-\bn{ছ}-} was, and the same scene is in the past:

\begin{quote}
\textbf{\bn{কাল} \bn{সন্ধ্যায়} \bn{আমি} \bn{বই} \bn{পড়ছিলাম।}}
\end{quote}

\begin{quote}
\emph{kāl sôndhyāy āmi bôi pôṛchhilām.} — \textquotedblleft{}Yesterday evening I was reading.\textquotedblright{}
\end{quote}

\begin{itemize}
\raggedright
  \item \textbf{\bn{আমি} \bn{পড়ছিলাম}} (\emph{āmi pôṛchhilām}) — I was reading
  \item \textbf{\bn{তুমি} \bn{পড়ছিলে}} (\emph{tumi pôṛchhile}) — you were reading
  \item \textbf{\bn{সে} \bn{পড়ছিল}} (\emph{she pôṛchhilô}) — he or she was reading
  \item \textbf{\bn{আপনি}} / \textbf{\bn{তিনি}} \textbf{\bn{পড়ছিলেন}} (\emph{pôṛchhilen}) — with respect
\end{itemize}

and \textbf{\bn{যাচ্ছিলাম}} (\emph{jāchchhilām}), \textquotedblleft{}I was going\textquotedblright{}. The plain past, \textbf{\bn{পড়লাম}}, is an event; \textbf{\bn{পড়ছিলাম}} is the scene around one.

\subsection*{Your turn}
Say these aloud:

\begin{itemize}
\raggedright
  \item \textquotedblleft{}She is eating\textquotedblright{} — \emph{she khāchchhe}
  \item \textquotedblleft{}We were going to the market\textquotedblright{} — \emph{āmrā bājāre jāchchhilām}
  \item \textquotedblleft{}What is happening?\textquotedblright{} — \emph{kī hôchchhe?}
\end{itemize}

\begin{tcolorbox}[breakable,colback=teal!4,colframe=teal!35!black,title={Before you move on}]
\textquotedblleft{}I'm going\textquotedblright{}? (\textbf{\bn{যাচ্ছি}}.) \textquotedblleft{}He was eating\textquotedblright{}? (\textbf{\bn{খাচ্ছিল}}.) Which is the scene and which the event: \textbf{\bn{পড়লাম}} or \textbf{\bn{পড়ছিলাম}}? (\textbf{\bn{পড়ছিলাম}} is the scene.)
\end{tcolorbox}
\section[\texorpdfstring{\bn{যখন} … \bn{তখন} … (jôkhon … tôkhon …)}{যখন … তখন … (jôkhon … tôkhon …)}]{\bn{যখন} … \bn{তখন} … (jôkhon … tôkhon …) — the scene and the event}

\label{lesson:BN-C266-jokhon-tokhon}



\begin{quote}
Say \textquotedblleft{}I was reading\textquotedblright{}. (\emph{āmi pôṛchhilām}.) You know \textbf{\bn{যখন}}, \textquotedblleft{}when\textquotedblright{}, and its partner \textbf{\bn{তখন}}, \textquotedblleft{}then\textquotedblright{}: \textbf{\bn{যখন} \bn{আমি} \bn{আসি}, \bn{তখন} \bn{চা} \bn{খাই}}. Read it aloud.
\end{quote}

\subsection*{You'll want to know: the scene and the event}
\begin{quote}
\textbf{\bn{যখন} \bn{আমি} \bn{বই} \bn{পড়ছিলাম}, \bn{তখন} \bn{মীরা} \bn{এল।}}
\end{quote}

\begin{quote}
\emph{jôkhon āmi bôi pôṛchhilām, tôkhon mīrā elô.} — \textquotedblleft{}I was reading when Mira came.\textquotedblright{}
\end{quote}

The \textbf{\bn{যখন}} half is the scene, in the going-on past. The \textbf{\bn{তখন}} half is the event that broke into it, in the plain past. English puts \textquotedblleft{}when\textquotedblright{} on the event; Bengali puts \textbf{\bn{যখন}} on the scene, and \textbf{\bn{তখন}} opens the event.

\subsection*{You'll want to know: \bn{তখন} on its own}
\textbf{\bn{তখন}} also stands alone in a story: \textquotedblleft{}then, at that time\textquotedblright{}.

\begin{quote}
\textbf{\bn{তখন} \bn{রাত} \bn{ছিল।} \bn{বাইরে} \bn{ঝড়} \bn{হচ্ছিল।}}
\end{quote}

\begin{quote}
\emph{tôkhon rāt chhilô. bāire jhôṛ hôchchhilô.} — \textquotedblleft{}It was night then. A storm was going on outside.\textquotedblright{}
\end{quote}

\subsection*{Your turn}
Say these aloud:

\begin{itemize}
\raggedright
  \item \textquotedblleft{}I was reading when Mira came\textquotedblright{} — \emph{jôkhon āmi bôi pôṛchhilām, tôkhon mīrā elô}
  \item \textquotedblleft{}It was night then\textquotedblright{} — \emph{tôkhon rāt chhilô}
  \item \textquotedblleft{}We were going to the market when the storm came\textquotedblright{} — \emph{jôkhon āmrā bājāre jāchchhilām, tôkhon jhôṛ elô}
\end{itemize}

\begin{tcolorbox}[breakable,colback=teal!4,colframe=teal!35!black,title={Before you move on}]
Which half takes \textbf{-\bn{ছিল}-}? (The \textbf{\bn{যখন}} half, the scene.) What does \textbf{\bn{তখন}} mean on its own? (\textquotedblleft{}Then, at that time\textquotedblright{}.)
\end{tcolorbox}
\section[\texorpdfstring{\bn{ঠিক} (ṭhik)}{ঠিক (ṭhik)}]{\bn{ঠিক} (ṭhik) — correct; exactly}

\label{lesson:BN-C266-thik}



\begin{quote}
Before the new one: say \textquotedblleft{}It was night then\textquotedblright{}. (\emph{tôkhon rāt chhilô}.)
\end{quote}

\subsection*{You'll want to know: \bn{ঠিক}}
\textbf{\bn{ঠিক}} — \emph{ṭhik} — \textquotedblleft{}correct; exactly\textquotedblright{}.

\textbf{\bn{তুমি} \bn{ঠিক}} (\emph{tumi ṭhik}) is \textquotedblleft{}you're right\textquotedblright{}, and \textbf{\bn{ঠিক} \bn{আছে}} (\emph{ṭhik āchhe}), \textquotedblleft{}it's all right\textquotedblright{}, is the everyday \textquotedblleft{}fine, agreed\textquotedblright{}. Before a time or a place, \textbf{\bn{ঠিক}} is \textquotedblleft{}exactly\textquotedblright{}:

\begin{quote}
\textbf{\bn{ঠিক} \bn{তখন} \bn{ট্রেন} \bn{এল।}}
\end{quote}

\begin{quote}
\emph{ṭhik tôkhon ṭren elô.} — \textquotedblleft{}Right then the train came.\textquotedblright{}
\end{quote}

\begin{cousinweb}
\textbf{\bn{ঠিক}} (\emph{ṭhik}) goes back to a form \emph{*ṭhīkka}, \textquotedblleft{}firm, right\textquotedblright{} (CDIAL 5503), which no text records (hence the star), and which CDIAL derives, probably, from Sanskrit \emph{sthita}, \textquotedblleft{}standing\textquotedblright{}. CDIAL lists the Bengali \textbf{\bn{ঠিক}}, \textquotedblleft{}correct\textquotedblright{}, under it.
\end{cousinweb}

\subsection*{Your turn}
Say these aloud:

\begin{itemize}
\raggedright
  \item \emph{ṭhik}
  \item \textquotedblleft{}You're right\textquotedblright{} — \emph{tumi ṭhik}
  \item \textquotedblleft{}Right then the train came\textquotedblright{} — \emph{ṭhik tôkhon ṭren elô}
\end{itemize}

\begin{tcolorbox}[breakable,colback=teal!4,colframe=teal!35!black,title={Before you move on}]
What does \textbf{\bn{ঠিক}} mean? (\textquotedblleft{}Correct\textquotedblright{}, and before a time, \textquotedblleft{}exactly\textquotedblright{}.) \textquotedblleft{}Fine, agreed\textquotedblright{}? (\textbf{\bn{ঠিক} \bn{আছে}}.)
\end{tcolorbox}
\section[\texorpdfstring{\bn{ঠ} (ṭhô)}{ঠ (ṭhô)}]{\bn{ঠ} (ṭhô) — the letter at the start of \bn{ঠিক}}

\label{lesson:BN-W266-letter-ttha}



\begin{quote}
Before the new one: say \textbf{\bn{ঠিক}}, \textquotedblleft{}correct\textquotedblright{}. Then read \textbf{\bn{ট}} aloud (\emph{ṭô}), the letter of \textbf{\bn{টাকা}}.
\end{quote}

\subsection*{Script you'll notice: \bn{ঠ}}
\textbf{\bn{ঠ}} — \emph{ṭhô}: the \textbf{\bn{ট}} sound with a breath after it, as \textbf{\bn{খ}} is \textbf{\bn{ক}} with a breath. It opens \textbf{\bn{ঠিক}}.

In the print of this book, \textbf{\bn{ঠ}} has a headline crossed by one curved line. Above the headline, the line rises in a curve that bows out to the left, its tip at the top. Below the headline it makes a large closed loop, shaped like a drop hanging from the headline, that fills the lower part of the letter; at its left side the loop has a sharp corner. \textbf{\bn{ঠ}} has no stem.

Set it beside \textbf{\bn{ট}}: \textbf{\bn{ট}} has a raised hook above its headline too, but below it a stem, and a bowl that ends in a solid knob. \textbf{\bn{ঠ}} has no knob.

\subsection*{Writing: \bn{ঠ} — copy the printed shape}
Put your pen beside \textbf{\bn{ঠ}} and copy the shape you can see, slowly and large. Check the copy against the print: the curve above the headline, the line crossing it, the closed loop and its corner.

\begin{quote}
This book does not yet tell you where to start \bn{ঠ} or which way the pen goes. That is not written down here until it can be written down with a source. Copying the printed shape needs no such source, and it is how the shape becomes yours in the meantime.
\end{quote}

\subsection*{Your turn}
\begin{itemize}
\raggedright
  \item \emph{Find:} \bn{ঠ} at the front of \textbf{\bn{ঠিক}}
  \item \emph{Write it:} \bn{ট}, then \bn{ঠ}, from the print, naming each
  \item \emph{Write it:} \textbf{\bn{ঠিক}}, then \textbf{\bn{ঠিক} \bn{আছে}}
  \item \emph{Write it:} \bn{ষ}, from memory
\end{itemize}

\begin{tcolorbox}[breakable,colback=teal!4,colframe=teal!35!black,title={Before you move on}]
Which letter opens \textbf{\bn{ঠিক}}? (\textbf{\bn{ঠ}}.) What does \textbf{\bn{ট}} have that \textbf{\bn{ঠ}} does not? (A stem and a solid knob.)
\end{tcolorbox}
\section[\texorpdfstring{\bn{ওঠা} (oṭhā)}{ওঠা (oṭhā)}]{\bn{ওঠা} (oṭhā) — to rise; to get up}

\label{lesson:BN-C266-otha}



\begin{quote}
Before the new one: write \textbf{\bn{ঠ}}, the letter of the last lesson, and say \textbf{\bn{ঠিক}}. The next word carries \textbf{\bn{ঠ}} too.
\end{quote}

\subsection*{You'll want to know: \bn{ওঠা}}
\textbf{\bn{ওঠা}} — \emph{oṭhā} — \textquotedblleft{}to rise; to get up\textquotedblright{}.

\begin{quote}
\textbf{\bn{সে} \bn{রোজ} \bn{সকালে} \bn{তাড়াতাড়ি} \bn{ওঠে।}}
\end{quote}

\begin{quote}
\emph{she roj sôkāle tāṛātāṛi oṭhe.} — \textquotedblleft{}She gets up early every morning.\textquotedblright{}
\end{quote}

\textbf{\bn{ওঠে}} is the present for he or she, and the moon and the sun \textquotedblleft{}rise\textquotedblright{} with it too: \textbf{\bn{রাতে} \bn{চাঁদ} \bn{ওঠে}} (\emph{rāte chā̃d oṭhe}), \textquotedblleft{}the moon rises at night\textquotedblright{}. In most other forms the \emph{o} rises to \emph{u}, as in \textbf{\bn{খোলা}} $\to$ \textbf{\bn{খুললাম}}: the past \emph{uṭhlô}, \textquotedblleft{}got up\textquotedblright{}, and the \emph{-e} form \emph{uṭhe}. At the start of a word that \emph{u} needs a letter of its own, which comes three chapters on; until then, say those forms, and write \textbf{\bn{ওঠে}}.

\begin{cousinweb}
\textbf{\bn{ওঠা}} (\emph{oṭhā}) goes back to a reconstructed \emph{*ut-sthāti}, \textquotedblleft{}stands up\textquotedblright{} (CDIAL 1900): \emph{ut-}, \textquotedblleft{}up\textquotedblright{}, on the root \emph{sthā}, \textquotedblleft{}to stand\textquotedblright{}. CDIAL lists the Bengali verb under it, as \emph{oṭhā} and \emph{uṭhā}, \textquotedblleft{}to get up, stand\textquotedblright{}.
\end{cousinweb}

\subsection*{Your turn}
Say these aloud:

\begin{itemize}
\raggedright
  \item \emph{oṭhā}, then \emph{she oṭhe}
  \item \textquotedblleft{}She gets up early every morning\textquotedblright{} — \emph{she roj sôkāle tāṛātāṛi oṭhe}
  \item \textquotedblleft{}The moon rises at night\textquotedblright{} — \emph{rāte chā̃d oṭhe}
\end{itemize}

\begin{tcolorbox}[breakable,colback=teal!4,colframe=teal!35!black,title={Before you move on}]
What does \textbf{\bn{ওঠা}} mean? (\textquotedblleft{}To rise; to get up\textquotedblright{}.) \textquotedblleft{}She gets up\textquotedblright{}? (\textbf{\bn{সে} \bn{ওঠে}}.) What does its \emph{o} become in the past? (\emph{u}: \emph{uṭhlô}.)
\end{tcolorbox}
\section[(review)]{Review — the scene and the event}

\label{lesson:BN-R266-the-scene}



\begin{quote}
Say \textquotedblleft{}I'm reading\textquotedblright{} and \textquotedblleft{}I was reading\textquotedblright{}. (\emph{pôṛchhi}, \emph{pôṛchhilām}.)
\end{quote}

\subsection*{Your turn}
\begin{itemize}
\raggedright
  \item \emph{Say it:} What are you doing now? I'm eating.
  \item \emph{Say it:} We were sitting at the station when the train came.
  \item \emph{Say it:} It was evening then. A storm was going on.
  \item \emph{Say it:} You're right. Fine, I'll come tomorrow.
  \item \emph{Say it:} Mira gets up early every day.
  \item \emph{Say it:} Right then, Mira phoned.
  \item \emph{Write it:} \textbf{\bn{ঠিক}}, then \bn{ঘ} and \bn{ষ}
\end{itemize}

\begin{tcolorbox}[breakable,colback=teal!4,colframe=teal!35!black,title={Before you move on}]
Which form says \textquotedblleft{}going on now\textquotedblright{}? (\textbf{-\bn{ছ}-}: \textbf{\bn{পড়ছি}}.) Which half of \textbf{\bn{যখন} … \bn{তখন}} is the scene? (The \textbf{\bn{যখন}} half.) \textquotedblleft{}Fine, agreed\textquotedblright{}? (\textbf{\bn{ঠিক} \bn{আছে}}.) \textquotedblleft{}She gets up\textquotedblright{}? (\textbf{\bn{সে} \bn{ওঠে}}.)
\end{tcolorbox}
